\chapter{Statistical Mechanics and the Origins of Learning Machines}
\label{ch:statmech}\label{ch:background}

Statistical mechanics is the discipline that explains how probabilistic
laws emerge from deterministic dynamics. It is also the field in which the
neural network was invented. This chapter serves both facts, and its two
halves have distinct jobs. The first half builds the small set of physical
tools that the rest of the thesis uses (the Boltzmann distribution, the
partition function and free energy, and the Ising energy), with the
derivations carried out in full: the Boltzmann distribution is the
stationary law of Langevin dynamics (Chapter~\ref{ch:stochproc}), the free
energy is the variational principle behind the diffusion-model training
objective (Chapter~\ref{ch:diffusion}), and sampling by simulating a
physical process is statistical mechanics turned into an algorithm. The
second half is the part of the literature review that most histories of
machine learning skip: how spin glasses became the Hopfield network, how
the Hopfield network became the Boltzmann machine, and how the statistical
mechanics of these systems grew into the modern theory of learning. This
is the tradition, recognised by the 2024 Nobel Prize in Physics, in which
this thesis works. The physical presentation follows classical references
\citep{goldstein2002classical, landau1980statistical}.

\section{Hamiltonian mechanics and phase space}\label{sec:hamiltonian}

Consider a mechanical system with $n$ degrees of freedom, described by
generalised coordinates $q = (q_1, \dots, q_n)$ and conjugate momenta
$p = (p_1, \dots, p_n)$. Its dynamics is generated by a single scalar
function, the \emph{Hamiltonian} $H(q, p)$, which for the systems
considered here has the mechanical form
\begin{equation}
  H(q, p) \;=\; \underbrace{\sum_{i=1}^{n} \frac{p_i^2}{2m_i}}_{\text{kinetic}}
  \;+\; \underbrace{U(q_1, \dots, q_n)}_{\text{potential}},
  \label{eq:hamiltonian}
\end{equation}
and equals the total energy of the system. The state of the system at any
instant is the pair $(q, p)$, a point in the $2n$-dimensional \emph{phase
space} $\Gamma \subseteq \R^{2n}$. The evolution is given by
\emph{Hamilton's equations},
\begin{equation}
  \dot{q}_i \;=\; \frac{\partial H}{\partial p_i},
  \qquad
  \dot{p}_i \;=\; -\,\frac{\partial H}{\partial q_i},
  \qquad i = 1, \dots, n.
  \label{eq:hamilton}
\end{equation}
Two structural properties of \eqref{eq:hamilton} are the pillars on which
statistical mechanics rests.

\paragraph{Energy conservation.} Along any trajectory,
\begin{equation}
  \frac{\dd H}{\dd t}
  = \sum_{i=1}^n \left(
      \frac{\partial H}{\partial q_i}\,\dot q_i
      + \frac{\partial H}{\partial p_i}\,\dot p_i
    \right)
  = \sum_{i=1}^n \left(
      \frac{\partial H}{\partial q_i}\frac{\partial H}{\partial p_i}
      - \frac{\partial H}{\partial p_i}\frac{\partial H}{\partial q_i}
    \right)
  = 0 .
  \label{eq:energy-conservation}
\end{equation}
The motion is therefore confined to the \emph{energy shell}
$\Gamma_E = \{(q,p) : H(q,p) = E\}$, a $(2n-1)$-dimensional hypersurface.

\paragraph{Incompressibility of the phase flow.} The phase-space velocity
field
\begin{equation}
  v(q, p) \;=\; (\dot q, \dot p)
  \;=\; \left( \frac{\partial H}{\partial p},\, -\frac{\partial H}{\partial q} \right)
  \label{eq:phase-velocity}
\end{equation}
has identically vanishing divergence:
\begin{equation}
  \nabla \cdot v
  \;=\; \sum_{i=1}^n \left(
      \frac{\partial \dot q_i}{\partial q_i}
      + \frac{\partial \dot p_i}{\partial p_i}
    \right)
  \;=\; \sum_{i=1}^n \left(
      \frac{\partial^2 H}{\partial q_i \partial p_i}
      - \frac{\partial^2 H}{\partial p_i \partial q_i}
    \right)
  \;=\; 0 .
  \label{eq:incompressible}
\end{equation}
Hamiltonian flow behaves like an incompressible fluid in phase space: it
can stretch and fold a region, but never change its volume. This identity
is the seed of the probabilistic construction that follows.

\section{Liouville's theorem}\label{sec:liouville}

A single trajectory is useless for macroscopic physics: for a gas,
$n \sim 10^{23}$ and neither the initial condition nor the solution of
\eqref{eq:hamilton} is accessible. The move that founds statistical
mechanics is to describe our \emph{ignorance} of the microstate by a
probability density $\rho(q, p, t)$ on phase space:
$\rho(q,p,t)\,\dd q\,\dd p$ is the probability that the system is in the
infinitesimal cell around $(q,p)$ at time $t$. Liouville's theorem states
how this density evolves.

\begin{theorem}[Liouville]\label{thm:liouville}
Under the Hamiltonian flow \eqref{eq:hamilton}, the phase-space density
obeys
\begin{equation}
  \frac{\partial \rho}{\partial t}
  \;=\; -\,\{\rho, H\}
  \;=\; \sum_{i=1}^n \left(
    \frac{\partial H}{\partial q_i}\frac{\partial \rho}{\partial p_i}
    - \frac{\partial H}{\partial p_i}\frac{\partial \rho}{\partial q_i}
  \right),
  \label{eq:liouville}
\end{equation}
where $\{\cdot,\cdot\}$ is the Poisson bracket. Equivalently, the density
is constant along trajectories: $\dd \rho / \dd t = 0$.
\end{theorem}

\begin{toolbox}{Proof of Liouville's theorem, step by step}
The density is locally conserved (probability mass is neither created nor
destroyed, only transported by the flow), so it satisfies the continuity
equation familiar from fluid dynamics:
\begin{equation*}
  \frac{\partial \rho}{\partial t} + \nabla \cdot (\rho\, v) = 0,
  \qquad v = (\dot q, \dot p).
\end{equation*}
Expand the divergence of the product using
$\nabla \cdot (\rho v) = v \cdot \nabla \rho + \rho\, (\nabla \cdot v)$:
\begin{equation*}
  \frac{\partial \rho}{\partial t}
  + \underbrace{\sum_{i=1}^n \left(
      \dot q_i \frac{\partial \rho}{\partial q_i}
      + \dot p_i \frac{\partial \rho}{\partial p_i}\right)}_{v \cdot \nabla \rho}
  + \rho \underbrace{\sum_{i=1}^n \left(
      \frac{\partial \dot q_i}{\partial q_i}
      + \frac{\partial \dot p_i}{\partial p_i}\right)}_{\nabla \cdot v}
  = 0 .
\end{equation*}
By the incompressibility identity \eqref{eq:incompressible} the last term
vanishes. Substituting Hamilton's equations
$\dot q_i = \partial H/\partial p_i$, $\dot p_i = -\partial H/\partial q_i$
into the middle term:
\begin{equation*}
  \frac{\partial \rho}{\partial t}
  + \sum_{i=1}^n \left(
      \frac{\partial H}{\partial p_i}\frac{\partial \rho}{\partial q_i}
      - \frac{\partial H}{\partial q_i}\frac{\partial \rho}{\partial p_i}
    \right) = 0
  \quad\Longleftrightarrow\quad
  \frac{\partial \rho}{\partial t} = -\{\rho, H\}.
\end{equation*}
Finally, the total (convective) derivative along a trajectory is
\begin{equation*}
  \frac{\dd \rho}{\dd t}
  = \frac{\partial \rho}{\partial t} + v \cdot \nabla \rho
  = \frac{\partial \rho}{\partial t} + \{\rho, H\}
  = 0 ,
\end{equation*}
which is the second form of the statement: an observer riding a
phase-space trajectory sees a constant local density. \hfill$\square$
\end{toolbox}

Liouville's theorem has an immediate consequence: \emph{any density that
is a function of conserved quantities only is stationary}. If
$\rho(q,p) = f\big(H(q,p)\big)$ for some function $f$, then
$\{\rho, H\} = f'(H)\,\{H, H\} = 0$, so $\partial \rho / \partial t = 0$.
Equilibrium statistical mechanics is the study of which $f$ nature
chooses.

\section{From dynamics to ensembles}\label{sec:ensembles}

An \emph{ensemble} is a probability distribution over microstates that is
stationary under the dynamics and compatible with the macroscopic
constraints we impose (fixed energy, fixed temperature, and so on). Two
ensembles matter for this thesis.

\subsection{The microcanonical ensemble}

For a perfectly isolated system, the energy is fixed at $E$ and
Liouville's theorem suggests the simplest choice compatible with it: a
uniform density on the energy shell,
\begin{equation}
  \rho_{\mathrm{mc}}(q, p)
  \;=\; \frac{1}{\Omega(E)}\,
        \delta\big(H(q,p) - E\big),
  \qquad
  \Omega(E) = \int_\Gamma \delta\big(H(q,p) - E\big)\,\dd q\,\dd p ,
  \label{eq:microcanonical}
\end{equation}
where $\Omega(E)$ counts the microstates at energy $E$. The postulate
that all accessible microstates are equally likely (the \emph{equal a
priori probability} postulate) is motivated, not proved, by ergodicity:
for ergodic dynamics, time averages along a single trajectory converge to
averages over the shell, so the uniform measure is the one a long
experiment effectively samples. Boltzmann's entropy attaches a number to
this count,
\begin{equation}
  S(E) \;=\; k_B \log \Omega(E),
  \label{eq:boltzmann-entropy}
\end{equation}
with $k_B$ the Boltzmann constant. Entropy is the logarithm of the number
of microscopic ways a macroscopic state can be realised.

\subsection{The canonical ensemble and the Boltzmann distribution}
\label{sec:canonical}

Isolated systems are an idealisation; the physically ubiquitous situation
is a system exchanging energy with a much larger environment (a
\emph{heat bath}) at temperature $T$. The stationary distribution of the
system alone is then no longer uniform: it is the
\emph{Boltzmann--Gibbs distribution},
\begin{equation}
  \boxed{\;
  \rho_\beta(x) \;=\; \frac{1}{Z(\beta)}\, e^{-\beta H(x)},
  \qquad
  Z(\beta) = \int e^{-\beta H(x)}\,\dd x,
  \qquad
  \beta = \frac{1}{k_B T},
  \;}
  \label{eq:boltzmann}
\end{equation}
where $x = (q,p)$ collects the microstate and $Z(\beta)$, the
\emph{partition function}, normalises the density. Equation
\eqref{eq:boltzmann} can be derived in two independent ways, and both
derivations are instructive enough to be carried out in full: the first
(physical) traces the exponential to the entropy of the bath; the second
(information-theoretic) characterises it as the least-committal
distribution compatible with a fixed mean energy, a perspective due to
\citet{jaynes1957information} that returns, almost verbatim, in modern
machine learning.

\begin{toolbox}{The Boltzmann distribution from the heat bath}
Let the total system (system $S$ together with bath $B$) be isolated with
total energy $E_{\mathrm{tot}}$, and let the coupling be weak, so
$H_{\mathrm{tot}} = H_S + H_B$. Apply the microcanonical postulate to the
total system: the probability that $S$ is in a \emph{specific} microstate
$x$ with energy $H_S(x) = \epsilon$ is proportional to the number of bath
microstates carrying the remaining energy,
\begin{equation*}
  \Prob(x) \;\propto\; \Omega_B\big(E_{\mathrm{tot}} - \epsilon\big)
  \;=\; \exp\!\left[ \tfrac{1}{k_B} S_B\big(E_{\mathrm{tot}} - \epsilon\big) \right],
\end{equation*}
using the definition \eqref{eq:boltzmann-entropy} of the bath entropy.
Since the bath is huge, $\epsilon \ll E_{\mathrm{tot}}$, expand $S_B$ to
first order:
\begin{equation*}
  S_B(E_{\mathrm{tot}} - \epsilon)
  = S_B(E_{\mathrm{tot}})
  - \epsilon \,\frac{\partial S_B}{\partial E}\Big|_{E_{\mathrm{tot}}}
  + O\!\big(\epsilon^2 / E_{\mathrm{tot}}\big).
\end{equation*}
The derivative of the bath entropy with respect to its energy
\emph{defines} the bath temperature,
$\partial S_B / \partial E = 1/T$. Substituting,
\begin{equation*}
  \Prob(x) \;\propto\;
  \underbrace{e^{S_B(E_{\mathrm{tot}})/k_B}}_{\text{constant in } x}
  \cdot\, e^{-\epsilon/(k_B T)}
  \;\propto\; e^{-\beta H_S(x)} .
\end{equation*}
The quadratic remainder is suppressed by the bath's heat capacity and
vanishes in the thermodynamic limit. Normalising yields
\eqref{eq:boltzmann}. Note what happened: the exponential form is nothing
but the first-order Taylor expansion of the bath's entropy. The Boltzmann
factor is the price, in bath microstates, of lending energy $\epsilon$ to
the system. \hfill$\square$
\end{toolbox}

\begin{toolbox}[tb:maxent]{The Boltzmann distribution from maximum entropy}
Forget physics; ask, with \citet{jaynes1957information}: among all
densities $\rho$ with a given mean energy $\E_\rho[H] = U$, which one
assumes the least beyond the constraint? ``Least'' is measured by the
Gibbs--Shannon entropy $S[\rho] = -k_B \int \rho \log \rho$. Maximise
$S[\rho]$ subject to $\int \rho = 1$ and $\int \rho H = U$ with Lagrange
multipliers $\lambda_0, \lambda_1$:
\begin{equation*}
  \mathcal{L}[\rho]
  = -k_B \!\int\! \rho \log \rho \,\dd x
  \;-\; \lambda_0 \Big( \!\int\! \rho \,\dd x - 1 \Big)
  \;-\; \lambda_1 \Big( \!\int\! \rho H \,\dd x - U \Big).
\end{equation*}
Take the functional derivative and set it to zero pointwise:
\begin{equation*}
  \frac{\delta \mathcal{L}}{\delta \rho(x)}
  = -k_B\big(\log \rho(x) + 1\big) - \lambda_0 - \lambda_1 H(x) = 0
  \;\;\Longrightarrow\;\;
  \rho(x) = \exp\!\Big[-1 - \tfrac{\lambda_0}{k_B} - \tfrac{\lambda_1}{k_B} H(x)\Big].
\end{equation*}
The density is an exponential of the constrained quantity. This is
completely general: constraining $\E[\phi_j(x)]$ for features $\phi_j$
produces $\rho \propto \exp[-\sum_j \beta_j \phi_j(x)]$, the exponential
family. The multiplier $\lambda_0$ is fixed by normalisation and produces
$Z$; renaming $\lambda_1/k_B = \beta$ and identifying $\beta = 1/(k_B T)$
through thermodynamics recovers \eqref{eq:boltzmann} exactly.
Second-order check:
$\delta^2 \mathcal{L} / \delta \rho^2 = -k_B/\rho < 0$, so the stationary
point is the unique maximum (the entropy is strictly concave).
\hfill$\square$
\end{toolbox}

The second derivation deserves a remark. The Boltzmann distribution is
not a statement about molecules but about \emph{inference under
constraints}. Any time we posit a model of the form
$p(x) \propto e^{-E(x)}$, as we will for Markov random fields
(Chapter~\ref{ch:graphs}) and energy-based models
(Chapter~\ref{ch:ebm}), we are doing statistical mechanics, whether the
variables are spins, pixels, or the frames of a video.

\section{The partition function and free energy}\label{sec:free-energy}

The partition function $Z(\beta)$ in \eqref{eq:boltzmann} looks like a
mere normaliser; it is in fact the central computational object of the
theory, and the reason approximate inference exists as a field. Its
logarithm defines the \emph{Helmholtz free energy},
\begin{equation}
  F(\beta) \;=\; -\frac{1}{\beta} \log Z(\beta).
  \label{eq:free-energy}
\end{equation}

\begin{toolbox}{Everything is a derivative of $\log Z$}
Differentiate $\log Z$ with respect to $\beta$, carrying the derivative
through the integral:
\begin{equation*}
  -\frac{\partial \log Z}{\partial \beta}
  = -\frac{1}{Z}\frac{\partial}{\partial \beta}\!\int\! e^{-\beta H}
  = \frac{\int H\, e^{-\beta H}}{\int e^{-\beta H}}
  = \E_{\rho_\beta}[H] \;=\; U ,
\end{equation*}
the mean energy. Differentiate once more:
\begin{align*}
  \frac{\partial^2 \log Z}{\partial \beta^2}
  &= \frac{\partial}{\partial \beta}
     \left( -\frac{\int H e^{-\beta H}}{Z} \right)
   = \frac{\int H^2 e^{-\beta H}}{Z}
     - \left( \frac{\int H e^{-\beta H}}{Z} \right)^{\!2} \\[2pt]
  &= \E[H^2] - \E[H]^2 \;=\; \operatorname{Var}(H) \;\geq\; 0 .
\end{align*}
So $\log Z$ is convex in $\beta$, its first derivative gives the mean
energy, its second the energy fluctuations; the latter equals
$k_B T^2 C_V$ with $C_V$ the heat capacity. This is a
\emph{fluctuation--dissipation} relation: how a system responds to
heating is determined by its spontaneous equilibrium fluctuations.
Finally, rewriting the Gibbs--Shannon entropy of $\rho_\beta$:
\begin{equation*}
  S[\rho_\beta]
  = -k_B \E\big[\log \rho_\beta\big]
  = -k_B \E\big[-\beta H - \log Z\big]
  = \frac{U}{T} + k_B \log Z ,
\end{equation*}
and multiplying by $-T$: $\; -TS = -U + F$, i.e.
\begin{equation*}
  F = U - TS .
\end{equation*}
The free energy is the exact trade-off between energy and entropy: at low
$T$ minimising $F$ means minimising energy (ordered states win), at high
$T$ it means maximising entropy (disordered states win). Phase
transitions are the non-analyticities of $F$ where the winner changes.
\hfill$\square$
\end{toolbox}

There is one more characterisation of $F$ worth recording, because it is
the statistical-mechanics form of the variational principle that
underlies both the training objective of diffusion models
(Chapter~\ref{ch:diffusion}) and the Bethe free energies of message
passing. For \emph{any} trial density $\mu$,
\begin{equation}
  \underbrace{\E_\mu[H] - \tfrac{1}{\beta} S[\mu]/k_B}_{\text{variational free energy } F[\mu]}
  \;=\; F \;+\; \tfrac{1}{\beta}\, \kl{\mu}{\rho_\beta}
  \;\;\geq\;\; F ,
  \label{eq:gibbs-variational}
\end{equation}
with equality iff $\mu = \rho_\beta$: the Boltzmann distribution is the
unique minimiser of the variational free energy, and the gap is exactly a
Kullback--Leibler divergence. Every variational method in machine
learning is a restriction of \eqref{eq:gibbs-variational} to a tractable
family of trial densities.

\section{Interacting systems: the Ising paradigm}\label{sec:ising}

Everything so far holds for arbitrary $H$; the physics (and the
difficulty) enters when the energy couples many degrees of freedom. The
canonical example is the \emph{Ising model} \citep{ising1925beitrag}:
binary spins $\sigma_i \in \{-1, +1\}$ on the nodes of a graph
$G = (V, \mathcal{E})$, with energy
\begin{equation}
  H(\sigma) \;=\; -\sum_{(i,j) \in \mathcal{E}} J_{ij}\, \sigma_i \sigma_j
  \;-\; \sum_{i \in V} h_i\, \sigma_i ,
  \label{eq:ising}
\end{equation}
where $J_{ij}$ are pairwise couplings and $h_i$ external fields. The
Gibbs distribution $p(\sigma) \propto e^{-\beta H(\sigma)}$ makes
neighbouring spins prefer alignment when $J_{ij} > 0$. Three lessons from
this model shape the rest of the thesis.

\paragraph{(i) The partition function is the hard part.}
$Z = \sum_\sigma e^{-\beta H(\sigma)}$ runs over $2^{|V|}$
configurations. On two-dimensional lattices Onsager computed it exactly
\citep{onsager1944crystal}; on general graphs the computation is
\#P-hard, and all of approximate inference (Monte Carlo sampling,
variational methods, message passing) exists because of this wall. On
\emph{trees}, however, $Z$ and all marginals are computable exactly in
time linear in the system size by the transfer-matrix recursion, the same
recursion that Chapter~\ref{ch:graphs} presents as belief propagation.
This tractable island is precisely where our research model lives.

\paragraph{(ii) Collective order emerges.} For $\beta$ beyond a critical
value the two-dimensional model magnetises spontaneously: a macroscopic
fraction of spins align even at $h = 0$, and the symmetric state splits
into two ergodic components. This spontaneous symmetry breaking is the
archetype of a \emph{phase transition}, and its dynamical counterpart, a
generation trajectory committing to one mode of a multimodal
distribution, is exactly the \emph{speciation} phenomenon of diffusion
models \citep{biroli2024dynamical, raya2023spontaneous}.

\paragraph{(iii) Disorder creates complexity.} Making the couplings
$J_{ij}$ random (of both signs) yields spin glasses, and, as the next
sections recount, spin glasses yielded neural networks. This is where the
physics of this chapter stops being background and becomes the
literature this thesis belongs to.

\section{Disorder, frustration, and spin glasses}\label{sec:spinglass}

In a ferromagnet all couplings agree, and the two ground states (all
spins up, all spins down) are found by inspection. Draw the couplings
$J_{ij}$ at random with both signs \citep{edwards1975theory,
sherrington1975solvable} and this transparency collapses. A triangle with
two positive bonds and one negative bond already cannot satisfy all
three: whatever the spins do, at least one bond pays. This is
\emph{frustration}, and its large-$N$ consequence is a free-energy
landscape shattered into exponentially many valleys separated by
extensive barriers: the system has many near-equally-good ways of being,
none of them readable off the couplings. The mean-field theory of this
phenomenon, developed through the replica and cavity methods
\citep{mezard1987spin}, and the Thouless--Anderson--Palmer (TAP)
self-consistency equations \citep{thouless1977solution}, took a decade to
control and earned Parisi the 2021 Nobel Prize in Physics.

Two exports from this episode matter here. First, a language: rugged
landscapes, metastable states, order parameters that measure similarity
rather than magnetisation. Machine learning speaks this language every
time it discusses loss surfaces and local minima. Second, a technology:
the cavity and TAP equations are recursions on \emph{local} quantities
that solve a \emph{global} inference problem, and they are the direct
ancestors of the message-passing algorithms \citep{mezard2009information,
zdeborova2016statistical, krzakala2024statphys} that
Chapter~\ref{ch:graphs} develops and that the research chapters deploy on
diffusion scores. The path from disordered magnets to algorithms is not
an analogy; it is a genealogy.

\section{The Hopfield network: memory as a phase of matter}
\label{sec:hopfield}

\citet{hopfield1982neural} asked a physicist's question about memory: can
a system of simple binary units \emph{store} and \emph{retrieve} patterns
as a collective effect, the way a magnet stores a magnetisation? His
construction is an Ising system \eqref{eq:ising} in disguise. Take $N$
binary neurons $\sigma_i \in \{-1,+1\}$ and $P$ patterns
$\xi^\mu \in \{-1,+1\}^N$ to be memorised, and wire the couplings by the
\emph{Hebb rule}, ``neurons that fire together wire together''
\citep{hebb1949organization}:
\begin{equation}
  J_{ij} \;=\; \frac{1}{N} \sum_{\mu=1}^{P} \xi_i^\mu \xi_j^\mu ,
  \qquad
  H(\sigma) \;=\; -\frac{1}{2}\sum_{i \neq j} J_{ij}\, \sigma_i \sigma_j .
  \label{eq:hopfield}
\end{equation}
The dynamics is the simplest imaginable: pick a neuron, align it with its
local field, repeat,
\begin{equation}
  \sigma_i \;\leftarrow\; \operatorname{sign}\Big( \textstyle\sum_{j \neq i}
  J_{ij}\, \sigma_j \Big).
  \label{eq:hopfield-dynamics}
\end{equation}

\begin{toolbox}{The Hopfield dynamics descends the energy}
Let $h_i = \sum_{j \neq i} J_{ij} \sigma_j$ be the local field on neuron
$i$ and suppose the update \eqref{eq:hopfield-dynamics} flips it,
$\sigma_i' = -\sigma_i$. Only the terms of $H$ containing $\sigma_i$
change, so
\begin{equation*}
  \Delta H
  \;=\; H(\sigma') - H(\sigma)
  \;=\; -\,(\sigma_i' - \sigma_i)\, h_i
  \;=\; -\,2\,\sigma_i'\, h_i .
\end{equation*}
The update sets $\sigma_i' = \operatorname{sign}(h_i)$, hence
$\sigma_i' h_i = |h_i| \geq 0$ and $\Delta H \leq 0$, with strict
decrease whenever a flip actually occurs. Since $H$ is bounded below on a
finite configuration space, the dynamics must terminate in a local
minimum: \emph{every} initial condition falls into an attractor. Memory
retrieval is this fall: initialise at a corrupted pattern, descend, and
read off the fixed point. \hfill$\square$
\end{toolbox}

The patterns are, by construction, near the bottoms of the energy
\eqref{eq:hopfield}: retrieval works because each memory is an attractor
with a basin. How many memories fit is a genuine phase-transition
question, answered with spin-glass tools by \citet{amit1985storing}: for
random patterns with $P = \alpha N$, retrieval survives only below a
critical load $\alpha_c \approx 0.138$, beyond which the memories are
crushed into a useless glassy phase. Memory, in this model, is literally
a phase of matter, entered and left through sharp transitions. The
single-pattern case $P=1$, known as the Mattis model
\citep{mattis1976solvable}, is gauge-equivalent to the uniform
ferromagnet: flipping $\sigma_i \to \xi_i \sigma_i$ maps the stored
pattern onto the all-up state. This observation is worth keeping in mind:
the Mattis model returns in Section~\ref{sec:cascade} as the exactly
solvable teacher in a modern analysis of how energy-based models learn.

In 2024 the Nobel Prize in Physics went to John Hopfield and Geoffrey
Hinton ``for foundational discoveries and inventions that enable machine
learning with artificial neural networks'' \citep{nobel2024physics}:
formal recognition that the lineage traced in this section is physics,
not just inspired by it.

\section{From Hopfield to Boltzmann machines: learning enters}
\label{sec:boltzmann-machines}

The Hopfield network stores what its designer wires into it. The decisive
step towards machine \emph{learning} was to make the couplings themselves
the unknowns. Promote the zero-temperature dynamics
\eqref{eq:hopfield-dynamics} to stochastic updates at temperature
$T = 1/\beta$ (each neuron aligns with its field only with probability
given by the Boltzmann factor); the stationary law of this dynamics is
exactly the Gibbs measure $p(\sigma) \propto e^{-\beta H(\sigma)}$ of the
Ising energy. A \emph{Boltzmann machine} \citep{ackley1985learning} is
this object run backwards: fix a dataset, and ask for the couplings $J$
that make the Gibbs measure reproduce it. Gradient ascent on the
log-likelihood gives a rule of striking symmetry,
\begin{equation}
  \frac{\partial}{\partial J_{ij}} \,\E_{\text{data}}\big[\log p(\sigma)\big]
  \;=\; \beta\,\Big(
  \underbrace{\langle \sigma_i \sigma_j \rangle_{\text{data}}}_{\text{positive phase}}
  \;-\;
  \underbrace{\langle \sigma_i \sigma_j \rangle_{\text{model}}}_{\text{negative phase}}
  \Big),
  \label{eq:bm-learning}
\end{equation}
and learning stops when the model's correlations match the data's. This
is the maximum-entropy principle of Toolbox~\ref{tb:maxent} turned into
an algorithm: matching prescribed moments with an exponential-family
measure. The rule also contains, in miniature, the central obstruction of
the field. The positive phase is an average over data, which is cheap.
The negative phase is an average over the \emph{model's} Gibbs measure: a
partition-function computation, the wall of Section~\ref{sec:ising},
approachable only through sampling
(Section~\ref{sec:mcmc}).

The \emph{restricted} Boltzmann machine (RBM) tames this cost by
architecture \citep{smolensky1986information}: split the units into a
visible layer $v$ and a hidden layer $h$ with couplings only across
layers, $E(v,h) = -v^\T W h - b^\T v - c^\T h$. On this bipartite graph
each layer is conditionally independent given the other, so Gibbs
sampling alternates two exact, fully parallel steps, and the
contrastive-divergence approximation to \eqref{eq:bm-learning} made
training practical \citep{hinton2002training}. Stacked RBMs, the deep
belief networks, delivered the first working deep architectures and
relaunched neural networks as a field \citep{hinton2006fast}. The lineage
then split. The supervised branch, powered by backpropagation
\citep{rumelhart1986learning}, scaled through convolutional networks to
the modern deep-learning era \citep{krizhevsky2012imagenet,
lecun2015deep} and largely left equilibrium physics behind. The
\emph{generative} branch, the one this thesis lives in, kept the Gibbsian
core: a model is an energy, learning is moment matching, sampling is
dynamics. Chapter~\ref{ch:diffusion} follows that branch from
energy-based models to diffusion.

\section{The statistical mechanics of learning}\label{sec:statmech-learning}

Physics did not stop contributing machines; it also built the theory of
their behaviour. \citet{gardner1988space} showed that learning itself can
be treated as an equilibrium problem: fix a set of examples and compute,
by replica methods, the volume of coupling space compatible with them;
capacity, generalisation, and their transitions then come out as
thermodynamics in \emph{weight} space. This programme, extended through
teacher--student settings, matured into the modern statistical physics of
inference \citep{zdeborova2016statistical}, whose central lesson recurs
throughout this thesis: high-dimensional inference problems have phases,
and algorithms fail or succeed at sharp boundaries.

Two strands of this literature frame our research questions directly.
The first is algorithmic: \citet{mezard2017meanfield} derived the
mean-field message-passing (TAP) equations of the Hopfield model and its
generalisations, exhibiting the equivalence between Hopfield networks and
RBMs and showing that retrieval can be organised as a message-passing
computation. It is the same machinery, transplanted from memories to
diffusion scores, that the research chapters run on our model. The second
is dynamical: training itself undergoes phase transitions. In linear
networks, gradient descent learns the singular modes of the data
covariance sequentially \citep{saxe2014exact}; in RBMs, the weight
matrix's singular vectors align with the data's principal components in
the early phase of learning \citep{decelle2017spectral}, and hidden units
develop compositional receptive fields through condensation transitions
\citep{tubiana2017emergence}. The sharpest result in this line,
\citet{bachtis2024cascade}, examined in depth in
Section~\ref{sec:cascade}, shows that RBM training is an entire
\emph{cascade} of second-order phase transitions, one per learned
feature. Learning, like memory, turns out to be a sequence of symmetry
breakings.

\section{What this thesis takes from physics}

A \emph{complex system} is, operationally, a system whose macroscopic
behaviour is not readable off its microscopic rules: interactions
generate collective phenomena (order, phases, transitions, slow dynamics)
that must be derived, not postulated. Statistical mechanics is the
toolbox for that derivation, and after this chapter's story its objects
map one-to-one onto the objects of modern generative modelling:
\begin{center}
\begin{tabular}{ll}
\toprule
\textbf{Statistical mechanics} & \textbf{Generative modelling} \\
\midrule
microstate $x$ & data point / latent configuration \\
energy $H(x)$ & negative log-density $-\log p(x)$ (up to $Z$) \\
Boltzmann distribution $e^{-\beta H}/Z$ & model distribution \\
partition function $Z$ & normalising constant / evidence \\
free energy $F$ & negative ELBO (up to constants) \\
Hebbian couplings, attractors & stored memories, retrieval \\
moment matching \eqref{eq:bm-learning} & maximum-likelihood training \\
Langevin dynamics & sampling / reverse diffusion \\
phase transition & feature condensation, speciation \\
cavity / TAP equations & belief propagation / AMP \\
\bottomrule
\end{tabular}
\end{center}
The next chapter supplies the probabilistic dynamics, Markov chains and
their continuous-time limits, that turn the static ensembles of this
chapter into processes that move probability around, which is what a
diffusion model ultimately is.
